\section{Comparing the offsets of neighboring sheets}
\label{sec:neighbor-offsets}

Proposition~\ref{prop:neighbor-labels} gives a common linear label for
neighboring sheet profiles.  It does not yet identify their constant terms.
At a fixed positive chart enlargement, the direct bridge between the sheets
can have length proportional to that enlargement times the scale.  We retain
this error and compare volumes in suitable nilpotent quotients.  Their Haar
multipliers will detect each component of the difference of the offsets.

\subsection{Quotients associated to cones of weights}

We first construct the quotients used in the comparison.  A word in the
generators $X_i$ is called \emph{pure for $S$} if every letter has weight in
$S$, and \emph{mixed for $S$} otherwise.  These terms will always refer to the
cone specified in the same argument.

\begin{lemma}[Cone quotients and their Haar multipliers]
\label{lem:cone-quotients}
Let $S\subset W^*$ be either the closed positive ray generated by a generator
weight $\alpha$, or the closed cone generated by two linearly independent
generator weights $\alpha,\lambda$.  The span $\mathfrak i_S$ of the mixed
words is a $\mathfrak d$-invariant ideal of $\mathfrak n$ containing
$\mathfrak d\cap\mathfrak n$.  Put
\[
 I_S=\exp_N(\mathfrak i_S),\qquad L_S=N/I_S.
\]
The group $I_S$ is closed, and $L_S$ is connected, simply connected, and
nilpotent.  Every compact subset of $L_S$ has compact lifts in $N$ and hence
in $G$.  There is a well-defined action of $G$ on $L_S$ under which $N$ acts
by left translations and $D$ acts by its induced conjugation automorphisms.  Write
\[
 p_S(g)=g\cdot 1_{L_S}.
\]
Then $p_S(gh)=g\cdot p_S(h)$ and the stabilizer of $1_{L_S}$ is $I_SD$.
There is a linear form $\ell_S\in W^*$ such that, for every Haar-measurable
$E\subset L_S$,
\begin{equation}
 \vol_{L_S}(g\cdot E)
   =e^{\ell_S(\pi g)}\vol_{L_S}(E).
 \label{eq:cone-haar-multiplier}
\end{equation}
The form $\ell_S$ is the sum of the weights of
$\mathfrak n/\mathfrak i_S$, counted with their dimensions.  It can be
expressed as a nonnegative integral combination of the inside generator
weights, with a positive coefficient for every inside generator.  In
particular, it is a positive multiple of $\alpha$ in the ray case, and it has
the form $c\alpha+d\lambda$, with $c,d>0$, in the two-generator case.
\end{lemma}

\begin{proof}
Bracketing a mixed word with any generator produces another mixed word.
The derivation identity and the Jacobi identity therefore show that their
span is the ideal generated by the outside letters.  Homogeneous bracket
words span $\mathfrak n$.  If a weight $\gamma$ lies outside $S$, a word of
weight $\gamma$ cannot be pure: a sum of inside weights belongs to $S$.
Consequently the entire generalized weight space
$\mathfrak n_\gamma$ lies in $\mathfrak i_S$ whenever $\gamma\notin S$.
For $v\in\mathfrak d$ and an outside letter $X_i$, the vector $[v,X_i]$
stays in $\mathfrak n_{\alpha_i}$ and hence in $\mathfrak i_S$.
Applying the derivation identity to words proves
$[\mathfrak d,\mathfrak i_S]\subset\mathfrak i_S$.  This argument includes
the nilpotent parts of the generalized weight actions.

Both possible cones $S$ are pointed.  A nonempty sum of their nonzero
inside weights cannot be zero.  The quotient algebra is spanned by pure
words, so it has no zero generalized weight space.  The zero weight space
of $\mathfrak n$ is therefore contained in $\mathfrak i_S$; in particular
$\mathfrak d\cap\mathfrak n\subset\mathfrak i_S$.
The exponential descriptions of the subgroups give $N\cap D\subset I_S$.
The ideal exponentiates to a closed subgroup because the exponential map
is a diffeomorphism and $\mathfrak i_S$ is a closed linear subspace.
The Baker--Campbell--Hausdorff group on $\mathfrak n/\mathfrak i_S$
identifies with $N/I_S$: the algebra quotient induces a surjective group
homomorphism with kernel $I_S$.  This realizes $L_S$ as a connected,
simply connected nilpotent Lie group.  A linear section of the algebra
quotient, composed with its logarithm and the exponential map of $N$,
gives a continuous section $L_S\to N$.  Images of compact sets under
this section are compact, proving the lifting assertion.

The action of $N\rtimes D$ on $L_S$ is
\[
 (n,d)\cdot(xI_S)=n(dxd^{-1})I_S.
\]
The kernel of the multiplication map $N\rtimes D\to G$ consists of
$(h^{-1},h)$ with $h\in N\cap D$.  Its action sends $xI_S$ to
$xh^{-1}I_S=xI_S$.  Thus the action factors through $G$, without requiring
that the multiplication map be injective.  If $g=nd$, its value at the
identity is $nI_S$, proving the stabilizer assertion and the stated
equivariance of $p_S$.

Translations by $N$ preserve quotient Haar measure.  The determinant of the
$D$-action is the exponential of its trace character; its unipotent parts
have determinant one.  The weight description thus gives
\eqref{eq:cone-haar-multiplier}.  To check its positivity assertion, project
to the abelianization of $\mathfrak n$.  The image of $\mathfrak i_S$ there
is exactly the span of the outside generator images.  The inside generator
images remain linearly independent in the quotient abelianization.  They
can be included in a basis of $\mathfrak n/\mathfrak i_S$ selected from pure
words.  Summing that basis's weights proves the asserted integral
combination.  Every inside weight in a two-generator cone is a nonnegative
real combination of its two defining weights, and each defining generator
itself survives.  Both extremal coefficients are therefore positive.
\end{proof}

For a ray we also use the notation $I_\alpha,L_\alpha,p_\alpha,\ell_\alpha$.
These objects depend on the oriented ray, rather than its representative.
For $A\in\mathcal A$ and a cone $S$, put
$S'=(A^*)^{-1}S$.  Lemma~\ref{lem:support-symmetries} ensures that $S'$ is a
cone of the same permitted kind, defined by the corresponding target
generator rays.

\subsection{Ideal boxes and a covering with controlled cardinality}

The quotient Haar multiplier will measure the difference between the two
offsets.  To make that comparison, we will cover the ideal coordinates by
translates of a smaller box and bound their number by the ratio of the
two box volumes.

Fix a cone $S$ as above and a Euclidean norm on $W$.  Write $I:=I_S$.
Define
\[
 T_* = \{v\in W:\|v\|\le1,\ \beta(v)\le0\text{ for every }\beta\in S\},
 \qquad T_j=\epsilon_jT_*,\quad j\in\{2,1,+\},
\]
where $0<\epsilon_2<\epsilon_1<\epsilon_+$.  An inside weight has support
maximum zero on $T_*$.  If $\beta\notin S$, separation from the closed
convex cone gives a vector nonpositive on $S$ and strictly positive on
$\beta$.  Rescaling that vector into the unit ball proves
$\max_{T_*}\beta>0$.  Hence
\begin{equation}
 \begin{aligned}
 h_Z(T_j)&=\epsilon_j c_Z,\qquad c_Z:=h_Z(T_*)>0
       &&\text{for every mixed word }Z,\\
 h_Z(T_j)&=0&&\text{for every pure word }Z.
 \end{aligned}
 \label{eq:cone-word-widths}
\end{equation}

Select a maximal-width basis of $\mathfrak i_S$ from mixed words and use it
to define the ideal box $B^I_r(T_j,L)$, by the same scaled-log-coordinate
construction as $B_r(T_j,L)$.  The same basis can be used for all three
$j$, since their widths differ by a positive scalar factor.  Complete it
to a basis of $\mathfrak n$ by pure words whose quotient images form a
basis.  This is a maximal-width basis of the whole algebra.  Indeed any
basis with too few mixed vectors can replace an excess pure vector by a
mixed vector, increasing its total width.  Once its mixed vectors form a
basis of $\mathfrak i_S$, their maximal possible total is the one selected
above, and the remaining widths are zero.  Thus, with $P_j=P(T_j)$,
\begin{equation}
 \vol_I B^I_r(T_j,L)\asymp L^{\dim I}e^{rP_j}
 \quad\text{for fixed }L>0.
 \label{eq:ideal-box-volume}
\end{equation}
The normalized multiplication and word-expansion bounds of
Lemma~\ref{lem:box-calculus} apply to these ideal boxes as well: brackets
of mixed words remain mixed, and their word widths add.

We will also use coordinates separating this ideal from its quotient.
Choose a linear section
$\sigma:\mathfrak n/\mathfrak i_S\to\mathfrak n$ whose image is spanned by
the selected pure words.  The map
\begin{equation}
 (i,z)\longmapsto i\exp_N\bigl(\sigma(\log_{L_S}z)\bigr)
 \label{eq:ideal-section-coordinates}
\end{equation}
is a bijection from $I_S\times L_S$ to $N$.  Projection determines $z$,
and then determines $i$.  Both this map and its inverse are polynomial in
log coordinates by the nilpotent product formula.  We verify explicitly
that product Haar measure in these coordinates is Haar measure on $N$,
up to a fixed scalar.  Write $s(z)=\exp_N(\sigma(\log_{L_S}z))$.
For $a\in N$, with quotient image $\bar a\in L_S$, left multiplication
acts on these coordinates by
\[
 a\,i\,s(z)=\bigl(\Ad_a(i)c_a(z)\bigr)s(\bar a z),
 \qquad c_a(z)=a s(z)s(\bar a z)^{-1}\in I.
\]
The inner automorphism $\Ad_a|_I$ has determinant one, since its
differential is the exponential of a nilpotent derivation.  The group
$I$ is unimodular, so right translation by $c_a(z)$ also preserves its
Haar measure.  Meanwhile $z\mapsto\bar a z$ preserves Haar measure on
$L_S$.  Fubini proves invariance of the pushed-forward product measure
under every left translation by $a$.  The coordinate diffeomorphism
makes this measure a nonzero Radon measure, so uniqueness of Haar
measure proves the assertion.
For a common sequential subexponential envelope, the box calculus gives
\begin{equation}
 n\in B_r(T_j,e^{o(r)})
 \quad\Longleftrightarrow\quad
 i\in B^I_r(T_j,e^{o(r)}),\quad
 \|\log_{L_S}z\|\le e^{o(r)},
 \label{eq:ideal-section-boxes}
\end{equation}
where the envelope may be enlarged on either side.  To see the forward
implication, use the adapted maximal basis: projection retains only its
zero-width pure coordinates.  Multiplying by the inverse section then
leaves an ideal element in the same enlarged box.  Conversely, the section
of a subexponential quotient log ball lies in the zero-width part of the
box, and multiplication gives the assertion.

\begin{lemma}[Strict-width covering and absorption]
\label{lem:strict-width-tiling}
Fix the cone, radii, and a letter $X_i$ with weight $\alpha\in S$.
Fix also $M>0$.  For all sufficiently small fixed $\delta>0$, the following
claims hold along every sequence $r\to\infty$.  Suppose
\[
 k_q=\exp(t_qX_i),\qquad |t_q|\le e^{M\delta r+o(r)}.
\]
There is a measurable partition of $B^I_r(T_1)$ into sets indexed by anchors
$i_\kappa\in B^I_r(T_1)$, with
\begin{align}
 \text{tile}_\kappa&\subset i_\kappa k_q B^I_r(T_2,O(1))k_q^{-1},
 &\#\{\kappa\}&\le e^{r(P_1-P_2)+o(r)},
 \label{eq:cone-tiles}\\
 k_q^{-1}i_\kappa k_q&\in B^I_r(T_+,O(1)).
 \label{eq:cone-anchor}
\end{align}
Moreover, if $u\in B_r(T_2,e^{o(r)})$ and $e_r\in N$ satisfies
$\|\log_N e_r\|\le e^{M\delta r+o(r)}$, the ideal coordinate of $ue_r$ in
\eqref{eq:ideal-section-coordinates} belongs to
$B^I_r(T_2,e^{o(r)})$.  These statements hold simultaneously for any fixed
finite collection of the cones, templates, and their images under
$\mathcal A$, after decreasing $\delta$.
\end{lemma}

\begin{proof}
If $Z$ is mixed, adjoining occurrences of $X_i$ does not change its width
on $T_*$, because $s_i(T_*)=0$.  In
\[
 \Ad_{k_q}Z
  =\sum_{p=0}^{m-1}\frac{t_q^p}{p!}(\ad X_i)^p Z,
\]
every nonzero term is mixed and has the same $c_Z$.  Its added exponential
cost is at most $pM\delta r+o(r)$.  Expanding a word in a maximal-width
ideal basis uses only basis vectors whose widths are at least its own,
by basis exchange.  Consequently the positive differences
\[
 (\epsilon_1-\epsilon_2)c_Z,
 \qquad (\epsilon_+-\epsilon_1)c_Z
\]
absorb all those costs if $\delta$ is sufficiently small.  There are only
finitely many words and $p<m$.  Their positive gaps also absorb the common
$o(r)$ envelope.  We obtain
\[
 k_qB^I_r(T_2,1)k_q^{-1}\subset B^I_r(T_1,O(1)),
 \qquad
 k_q^{-1}B^I_r(T_1,1)k_q\subset B^I_r(T_+,O(1)).
\]

Put $Q_r=k_qB^I_r(T_2,1)k_q^{-1}$.  Choose a maximal disjoint family
$i_\kappa Q_r$ with anchors in $B^I_r(T_1)$.  Such a family is finite: each set
lies in $B^I_r(T_1,O(1))$, by the first inclusion and normalized
multiplication.  Conjugation by $k_q$ preserves Haar on $I_S$, since its
derivative is the exponential of the nilpotent map
$\ad(t_qX_i)|_{\mathfrak i_S}$ and has determinant one.
Dividing volumes in \eqref{eq:ideal-box-volume} bounds the number of
anchors by $O(e^{r(P_1-P_2)})$.  Maximality implies that every translate
$iQ_r$, for $i\in B^I_r(T_1)$, meets some $i_\kappa Q_r$.  Hence
\[
 B^I_r(T_1)\subset\bigcup_\kappa i_\kappa Q_rQ_r^{-1}
  \subset\bigcup_\kappa i_\kappa k_qB^I_r(T_2,O(1))k_q^{-1}.
\]
Assign a point to the first set containing it.  This gives the measurable
partition, and the second inclusion above proves \eqref{eq:cone-anchor}.

For the last assertion write $u=i\exp_N\sigma(s)$, where
$i\in B^I_r(T_2,e^{o(r)})$ and $\|s\|\le e^{o(r)}$, using
\eqref{eq:ideal-section-boxes}.  Factor
\[
 \exp_N\sigma(s)e_r=i_{e_r}\exp_N\sigma(s_{e_r}).
\]
The polynomial coordinate maps show that
$\|\log i_{e_r}\|+\|s_{e_r}\|\le e^{M'\delta r+o(r)}$ for a fixed multiple
$M'$ of $M$.  Every ideal basis width on $T_2$ is strictly positive.
Decrease $\delta$ so that $M'\delta$ is less than all these widths.
Then $i_{e_r}\in B^I_r(T_2,O(1))$ for all sufficiently large $r$, and the ideal
coordinate $ii_{e_r}$ has the required bound.  Finitely many simultaneous
requirements have a positive common threshold for $\delta$.
\end{proof}

\subsection{Conditioned probes and the offset comparison}

We retain the profiles of Definition~\ref{def:chart-profiles}:
\[
 f_{r,x,n}(z)=\frac{\pi F(xn\Delta_rz)-\pi F(x)}r,
 \qquad a\text{ uniform in }B_r(\mathcal C),\quad n=ak_0,
\]
where $k_0\in B_r((1+\delta)\mathcal C,e^{o(r)})$ is independent of $a$.
All subexponential bounds below use one common vanishing envelope for the
finite constructions under consideration.  The maps and basepoints may
vary along the sequence.  We use
\[
 \mathcal U_\delta
  =\{z\in D_\infty:\pi z\in(1-c\delta)\operatorname{int}\mathcal C\}.
\]

\begin{proposition}[Equality of neighboring offsets]
\label{prop:neighbor-offsets}
There exists $\delta_2>0$ such that, for every fixed
$0<\delta\le\delta_2$, the following holds.  Let $r\to\infty$, let $F$ range
over the fixed family of smoothed quasi-isometries, let $x\in G$ vary, and
let $k_0$ be any allowed deterministic suffix.  For a generator letter
$X_i$, put
\[
 k=\exp(tX_i),\qquad |t|\le e^{(1+\delta)r s_i(\mathcal C)+o(r)}.
\]
Every pair in the support of any limiting joint law of
$(f_{r,x,ak_0},f_{r,x,ak_0k})$ has identical affine restrictions to
$\mathcal U_\delta$.  Equivalently, on every fixed compact subset of that
domain, the difference of these profiles tends to zero in probability.
The conclusion also holds for suffixes and steps sampled independently of
$a$, provided their bounds have a common vanishing envelope.
\end{proposition}

\begin{proof}
Take $\delta$ below the threshold in Proposition~\ref{prop:neighbor-labels}.
The two profiles in a limiting support then have the form
\[
 h_1(z)=A\pi z+b_1,\qquad h_2(z)=A\pi z+b_2
 \quad(z\in\mathcal U_\delta)
\]
with one $A\in\mathcal A$.  Fix such a pair and a subsequence whose joint
laws converge to a law containing it in its support.  It suffices to prove
$b_1=b_2$.  There are finitely many letters and labels, so these may be
fixed after passing to a subsequence.  Write $\alpha=\alpha_i$.

\paragraph{Choice of templates.}
We will use the ray cone of $\alpha$, and then each cone generated by
$\alpha$ and a generator weight independent of it.  Fix one of these
cones $S$, and put $S'=(A^*)^{-1}S$.  Let $p_0$ be the contact point of the
unit ball with the $\alpha$ facet of $\mathcal C$.  This point belongs to
the relative interior of that facet.  Choose a fixed $c'>c$ and put
\[
 w_q=(1-c'\delta)p_0.
\]
Choose fixed radii $0<\epsilon_2<\epsilon_1<\epsilon_+$ as above, small
enough that
\begin{equation}
 w_q+T_j\Subset(1-c\delta)\operatorname{int}\mathcal C
 \quad(j=2,1,+)
 \label{eq:offset-template-interior}
\end{equation}
for every sufficiently small fixed $\delta$.  Here $\Subset$ means compact
containment.  To verify the choice, the other oriented-ray facets have
positive slack at $p_0$; a sufficiently small fixed $\epsilon_+$ preserves
their inequalities.  On the $\alpha$ facet, $\alpha(T_j)\le0$, while
$c'>c$ gives a strict margin at $w_q$.  The radii can therefore be chosen
before $\delta$ tends to zero.  On the target templates $AT_j$ the inside
weights of $S'$ have width zero and all mixed words have positive width.
Also
\[
 P(AT_j)=P(T_j)=P_j.
\]

The comparison will balance three counts.  We will construct at least
$e^{rP_1-o(r)}$ separated source endpoints and assign them to at most
$e^{r(P_1-P_2)+o(r)}$ tiles.  For each tile we will prove an upper
bound
\[
 e^{r(P_2-\ell_{S'}(b_2-b_1))+o(r)}
\]
for the separated image endpoints.  The term $P_2$ comes from their
ideal coordinates; the remaining term is the Haar multiplier on the
quotient.  Summing over tiles cancels $P_2$ and will force
$\ell_{S'}(b_2-b_1)\le0$.  The construction below arranges these two
coordinate restrictions on the same endpoints.

Normalize the source at $w_q$ using the exact automorphism $S_{-rw_q}$, and
write $k_q=S_{-rw_q}k$.  Since $\alpha(p_0)=s_i(\mathcal C)$,
\[
 k_q=\exp(t_qX_i),\qquad
 |t_q|\le e^{(1+c')\delta r s_i(\mathcal C)+o(r)}.
\]
Use Lemma~\ref{lem:strict-width-tiling} to partition $B^I_r(T_1)$, with
anchors $i_\kappa$.  Choose a fixed unit cube $J$ in
$\mathfrak n/\mathfrak i_S$, and write $\sigma$ for the section used in
\eqref{eq:ideal-section-coordinates}.

\paragraph{Three paths chosen independently of the starting sheet.}
Sample $i_1$ uniformly in $B^I_r(T_1)$ and $s$ uniformly in $J$,
independently of the original random $a$.  Let $i_\kappa$ be the anchor assigned
to $i_1$, so the tile index $\kappa$ depends only on this probe.
For the moment $n=ak_0$ is still random.  Define
\begin{equation}
 n_1=nS_{rw_q}(i_1\exp_N\sigma(s)),\qquad
 n_{2,\kappa}=nS_{rw_q}(i_\kappa)k,\qquad
 n_2=nS_{rw_q}(i_1)k.
 \label{eq:offset-probes}
\end{equation}
We will count the endpoints $n_1$.  The auxiliary endpoint $n_2$
keeps the ideal probe $i_1$ and replaces the quotient factor by the
neighbor step $k$, while $n_{2,\kappa}$ gives one second-sheet anchor
for its tile.  Their design gives the exact cancellation
\begin{equation}
 n_1^{-1}n_2=S_{rw_q}\bigl(\exp_N(-\sigma(s))k_q\bigr).
 \label{eq:offset-probe-cancellation}
\end{equation}
This difference is independent of $i_1$.  A path from $n$ to $n_1$
will control the target quotient coordinate relative to the first
sheet.  A path from $n_{2,\kappa}$ to $n_2$, followed by the bridge
in \eqref{eq:offset-probe-cancellation}, will control the target ideal
coordinate of the same endpoint relative to its second-sheet anchor.

The following three sheet changes can be generated with uniformly bounded
numbers of generator switches:
\begin{align*}
 n&\longrightarrow n_1
       &&\text{using the template }w_q+T_1,\\
 nk&\longrightarrow n_{2,\kappa}
       &&\text{using the template }w_q+T_+,\\
 n_{2,\kappa}&\longrightarrow n_2
       &&\text{using the template }w_q+T_2.
\end{align*}
For the first, $i_1\exp_N\sigma(s)\in B_r(T_1,O(1))$ by
\eqref{eq:ideal-section-boxes}.  For the second, the relative normalized
increment is $k_q^{-1}i_\kappa k_q$, which satisfies
\eqref{eq:cone-anchor}.  For the third it is
$k_q^{-1}i_\kappa^{-1}i_1k_q\in B^I_r(T_2,O(1))$ by
\eqref{eq:cone-tiles}.  Lemma~\ref{lem:product-schemes}, followed by
$S_{rw_q}$, gives the stated generation with generator amplitudes at most
$e^{rs_j(w_q+T_l)+o(r)}$ for the indicated template $l$.

Choose these product representations measurably as functions of
$(i_1,s)$, using the measurable bounded product selection in
Lemma~\ref{lem:product-schemes}.  Its finite family of bounded schemes can
be padded by identity factors, and its Borel local inverse choices give
Borel parameters for every endpoint.  The anchors are already measurable.
All choices, including the product parameters and the paths below, are
independent of $a$.  The number of switches is bounded independently of
the probe parameters.  Every intermediate suffix is a fixed-length
product of the original allowed suffix and increments in the template
boxes, and, when needed, $k$.  The box calculus and
\eqref{eq:offset-template-interior} keep it in
$B_r((1+\delta)\mathcal C,e^{o(r)})$, with one envelope for all probes.

Choose a pivot $d_q\in D$ with $\pi d_q=rw_q$ and $|d_q|_D=O(r)$, for
example by exponentiating $r$ times a fixed linear lift of $w_q$ to
$\mathfrak d$.  Each generator switch can be performed by a $D$-leg from
the pivot to a support height of its template, the short switch there,
and a reversed $D$-leg on the new sheet.  The legs have length $O(r)$ and
the switches have length $o(r)$, by Lemmas~\ref{lem:log-distortion} and
\ref{lem:support-switch-paths}.  Their scaled $D$-coordinates lie in a fixed
compact subset of $\mathcal U_\delta$, by
Lemma~\ref{lem:d-scaling} and the strict containment
\eqref{eq:offset-template-interior}.  We fix a compact neighborhood of all
these coordinates for the remaining profile comparisons.

\paragraph{Why the intermediate profiles retain their starting offsets.}
This step uses only the affinity of individual limiting profiles.
Consider one of the selected generator switches, with template
$E\Subset(1-c\delta)\operatorname{int}\mathcal C$ and letter $X_j$.
There is a nonempty open set of heights in that domain on which
$\alpha_j(w)>s_j(E)$, because the support maximum is attained inside the
open domain.  On a smaller such open set the difference is bounded below
by a positive number.  At its $D$-positions, conjugating the switch back
gives exponentially decaying generator coefficients, up to polynomial
factors.  The distance of its two source sheet points, divided by $r$,
therefore tends to zero.  Their image profiles agree there in every
joint limit.  Proposition~\ref{prop:affine-profiles} makes each limiting
profile affine on $\mathcal U_\delta$, so equality on this open set gives
equality throughout that domain.

It follows by sequence compactness that the difference across each such
switch tends to zero in probability on the fixed compact set of legs.
This statement is uniform for the permitted parameter choices with their
common envelope: a violating sequence of choices would give a joint limit
contradicting the preceding open-set argument.  Conditioning on any
independent probe parameters and integrating proves the same assertion
for the measurable choices just made.  Since there are boundedly many
switches, with probability tending to one their intermediate profiles
all differ by $o(1)$ from the starting profile of their respective sheet.
The first path uses the profile of $n$; the other two use that of $nk$.

We now choose a single starting $n$ retaining a positive volume of probes.
For a fixed tolerance $\eta>0$, let $E_r(\eta)$ be the event that the
original pair of profiles is within $\eta$ of $(h_1,h_2)$ on the required
compact.  The support hypothesis gives
$\liminf_r\mathbb P(E_r(\eta))>0$.  Let $b_r(\eta)$ be the joint
probability, over the original chart sample and the independent probes,
that one of the intermediate profile comparisons has error exceeding
$\eta$.  The preceding paragraph gives $b_r(\eta)\to0$.
For large $r$ it is less than one quarter of
$\mathbb P(E_r(\eta))$.  Fubini's theorem then supplies an original sample
in $E_r(\eta)$ for which fewer than half of the probes fail these
comparisons.  First do this for each fixed tolerance and then let the
tolerance decrease sufficiently slowly with $r$.  We obtain a
deterministic starting $n=n_r$ and a measurable set of good probes of
probability at least $1/2$, on which all the required intermediate
profiles approach their designated affine maps.  This order of choices
does not condition on a null event, and it does not require every probe
to succeed simultaneously.

The generation selectors prescribe the paths; the probability law here
remains the original product of ideal Haar measure and Lebesgue measure
on $J$.  The preceding success events are measurable because the selected
parameters are Borel and the profiles depend continuously on their sheet
coordinates.  Inner regularity of this finite-dimensional product measure
allows us to retain a compact subset of good probes of probability at
least $1/4$.

The map $(i_1,s)\mapsto i_1\exp_N\sigma(s)$ is injective and carries
product Haar measure to Haar measure up to a fixed scalar.  The good
probe set thus has actual $N$-volume bounded below by a fixed positive
multiple of $e^{rP_1}$.  Since $S_{rw_q}$ has determinant one and left
translation by $n$ preserves Haar, the good values of $n_1$ in
\eqref{eq:offset-probes} have this same lower bound.  The possibly small
probability of the starting event has caused no loss of probe volume.
We use the compact retained subset from now on.  Its $n_1$-endpoint image
is compact, because $(i_1,s)\mapsto nS_{rw_q}(i_1\exp_N\sigma(s))$ is
continuous.  No continuity of the anchor assignment or generation
selectors is required; the tile partition is Borel and its anchor is
fixed on each tile.

\paragraph{The bridge and the two target restrictions.}
We now estimate the bridge in \eqref{eq:offset-probe-cancellation}.
At the pivot, conjugation by $d_q^{-1}$ cancels the character scaling
$S_{rw_q}$ and leaves only polynomially bounded unipotent factors.  As $s$
ranges in a fixed cube, the resulting logarithm has norm at most
$e^{C_1\delta r+o(r)}$.  Logarithmic distortion gives
\begin{equation}
 d_G(xn_1d_q,xn_2d_q)\le C_2\delta r+o(r),
 \label{eq:offset-bridge}
\end{equation}
where the constants are independent of sufficiently small fixed
$\delta$.  We retain this error at its stated size.

For good probes put
\[
 y_0=F(xnd_q),\quad y_1=F(xn_1d_q),\quad
 y_{2,\kappa}=F(xn_{2,\kappa}d_q),\quad y_2=F(xn_2d_q).
\]
The first probe path has target relative heights in
$r(AT_1+o(1))$, because its profiles approach $A\pi z+b_1$ and its pivot
height is $w_q$.  Its length is $O(r)$.  Path lifting and packing therefore
give a lift of $y_0^{-1}y_1$ with $N$-part in
$B_r(AT_1,e^{o(r)})$ and $D$-length $O(r)$.
Its endpoint height is $o(r)$, since both endpoint profiles approach
$Aw_q+b_1$.  Every pure quotient width is zero, so
\begin{equation}
 p_{S'}(y_0^{-1}y_1)\in \mathcal T_r,
 \qquad
 \mathcal T_r=\{z\in L_{S'}:\|\log z\|\le e^{o(r)}\},
 \qquad \vol_{L_{S'}}\mathcal T_r\le e^{o(r)}.
 \label{eq:offset-first-quotient}
\end{equation}
Choose its envelope large enough that the same restriction holds after
fixed bounded right thickening of the endpoints.  Indeed a bounded final
step keeps the first path's relative heights in $r(AT_1+o(1))$: its
unthickened endpoint height is $o(r)$ and $0\in AT_1$.  The same path
lifting estimate applies.

For a tile contributing a good probe, the second starting path and the
profile conditioning give
\begin{equation}
 \pi(y_0^{-1}y_{2,\kappa})=r(b_2-b_1)+o(r).
 \label{eq:offset-anchor-height}
\end{equation}
The path from $y_{2,\kappa}$ to $y_2$ has relative heights in
$r(AT_2+o(1))$, length $O(r)$, and endpoint height $o(r)$.
Choose its lift $u_2d_2$ with
\[
 u_2\in B_r(AT_2,e^{o(r)}),\qquad
 |d_2|_D=O(r),\qquad \pi d_2=o(r).
\]
A lift of the bridge from $y_2$ to $y_1$, using
\eqref{eq:offset-bridge} and the quasi-isometry bounds, has $N$ log norm
at most $e^{C_3\delta r+o(r)}$ and $D$-length $O(\delta r)+o(r)$.
Conjugating its $N$-part by $d_2$ changes this estimate only by a
subexponential factor.  We obtain a lift of $y_{2,\kappa}^{-1}y_1$ of the form
\begin{equation}
 u_2 e_r d,
 \qquad \|\log e_r\|\le e^{C_4\delta r+o(r)},\quad
 |d|_D=O(r),\quad \|\pi d\|\le O(\delta r)+o(r).
 \label{eq:offset-lift-with-error}
\end{equation}
The last height bound is only $O(\delta r)$, not $o(r)$.

Apply the absorption part of Lemma~\ref{lem:strict-width-tiling} to the
target ideal $I'=I_{S'}$ and template $AT_2$.  After decreasing $\delta$,
the ideal coordinate of $u_2e_r$ belongs to
$B^{I'}_r(AT_2,e^{o(r)})$.  More explicitly, factor
$u_2=i_2\exp\sigma'(s_2)$, with $\|s_2\|\le e^{o(r)}$.
Factoring $\exp\sigma'(s_2)e_r$ again produces an ideal error with log
norm $e^{O(\delta)r+o(r)}$.  Every mixed ideal width is strictly positive,
so this error fits in the ideal box.  This is an ideal-coordinate bound
uniform over the accompanying quotient coordinates; it does not require
their location to stay in a subexponential ball.

The bound remains valid under fixed bounded right thickening of the
lifts in \eqref{eq:offset-lift-with-error}.  Such thickening adds an
$N$-factor conjugated by $d$.  The stated $D$-length and
$O(\delta r)+o(r)$ height bounds make its logarithm at most
$e^{O(\delta)r+o(r)}$, and the same positive ideal widths absorb it.
This explicitly accounts for the non-sublinear height of the lift
relative to the second-sheet anchor.

All constants in the bridge and polynomial coordinate estimates can be
chosen independently of sufficiently small $\delta$: the radii are fixed,
the pivots stay in fixed compact height sets, and the quasi-isometry
constants are fixed.  There are finitely many cones, letters, and labels.
We choose $\delta_2$ small enough to satisfy all their strict gap and
absorption inequalities simultaneously.

The fixed-$\delta$ bridge has now been absorbed in the ideal-coordinate
bound.  It remains to compare quotient Haar volumes: their exact
multiplier will detect $\ell_{S'}(b_2-b_1)$.

\paragraph{Volume at an arbitrary translated quotient location.}
Fix a contributing tile and let $g_\kappa=y_{2,\kappa}^{-1}y_0$.  Equivariance in
Lemma~\ref{lem:cone-quotients}, together with
\eqref{eq:offset-first-quotient}, gives the exact restriction
\begin{equation}
 p_{S'}(y_{2,\kappa}^{-1}y_1)\in g_\kappa\cdot \mathcal T_r.
 \label{eq:offset-translated-quotient}
\end{equation}
It also holds for the thickened endpoints described above.  The Haar
multiplier and \eqref{eq:offset-anchor-height} imply
\begin{equation}
 \vol_{L_{S'}}(g_\kappa\cdot \mathcal T_r)
 \le e^{-r\ell_{S'}(b_2-b_1)+o(r)}.
 \label{eq:offset-quotient-volume}
\end{equation}
The set may be far from the quotient identity.  Its location has no
effect on this equality of Haar multipliers.

In the ideal and section coordinates \eqref{eq:ideal-section-coordinates},
every admissible $N$-part of the thickened lifts lies in a set with ideal
coordinate in $B^{I'}_r(AT_2,e^{o(r)})$ and quotient coordinate in
$g_\kappa\cdot \mathcal T_r$.  Quotient integration and
\eqref{eq:ideal-box-volume} bound its Haar volume by
\begin{equation}
 e^{r(P_2-\ell_{S'}(b_2-b_1))+o(r)}.
 \label{eq:offset-tile-N-volume}
\end{equation}
This is a bound by a containing product set.  It requires neither
independence of the coordinates nor a product description of the
admissible set, and holds at any translated quotient location.

To turn this into a separated-point bound, take image endpoints in this
tile with a fixed sufficiently large separation.  Choose the admissible
lifts relative to $y_{2,\kappa}$ and thicken each on the right by a fixed small
neighborhood in $N\rtimes D$.  Their projections lie in disjoint metric
balls, so these lifted neighborhoods are disjoint even if
$N\rtimes D\to G$ is not injective.  Each has the same positive product
Haar volume.  The thickened $D$-coordinates lie in an intrinsic $D$-ball
of radius $O(r)$, of polynomial volume.  The two coordinate restrictions
just proved still hold for all the neighborhoods.  Fubini's theorem and
\eqref{eq:offset-tile-N-volume} therefore bound their number by
\begin{equation}
 e^{r(P_2-\ell_{S'}(b_2-b_1))+o(r)}.
 \label{eq:offset-tile-count}
\end{equation}
The envelopes here are common to all contributing tiles, since the probe
paths, profile tolerances, and coordinate bounds were uniform over their
parameters.

The good $n_1$-volume was at least a constant times $e^{rP_1}$.
Lemma~\ref{lem:fixed-sheet-packing}, at the common pivot $d_q$, supplies
at least $e^{rP_1-o(r)}$ suitably separated source points from that set.
Take the separation large enough that their images are separated at the
constant used in \eqref{eq:offset-tile-count}.  Assign each chosen point
to its probe tile.  Summing \eqref{eq:offset-tile-count} over at most
$e^{r(P_1-P_2)+o(r)}$ tiles gives
\[
 e^{rP_1-o(r)}
 \le e^{r(P_1-\ell_{S'}(b_2-b_1))+o(r)}.
\]
Divide logarithms by $r$ and let $r\to\infty$.  We have proved
$\ell_{S'}(b_2-b_1)\le0$.

For the reverse inequality take the original starting suffix to be
$k_0k$ and the step to be $k^{-1}$.  The former is still in the allowed
enlarged box by normalized multiplication, with an enlarged common
subexponential envelope; the latter has the same generator amplitude
bound.  The swapped pair belongs to the swapped limiting support.
Repeating the argument gives $\ell_{S'}(b_1-b_2)\le0$.  Thus
\begin{equation}
 \ell_{S'}(b_2-b_1)=0
 \quad\text{for each of the specified cone tests.}
 \label{eq:offset-traces-vanish}
\end{equation}

\paragraph{The trace tests determine the offset.}
For the ray cone, its target trace is a positive multiple of
$\beta=(A^*)^{-1}\alpha$.  Hence $\beta(b_2-b_1)=0$.
For each generator weight $\lambda$ independent of $\alpha$, the target
two-generator trace is $c\beta+d\eta$, where
$\eta=(A^*)^{-1}\lambda$ and $c,d>0$.  Equation
\eqref{eq:offset-traces-vanish} then gives $\eta(b_2-b_1)=0$.
Weights proportional to $\alpha$ require no additional test.  Since the
generator weights span $W^*$, their target transforms do as well, proving
$b_2=b_1$.  When $\dim W=1$, the ray test alone gives this conclusion.

Finally, compactness converts the support assertion into convergence in
probability on each compact subset: any violation would admit a joint
limit assigning positive mass to unequal restrictions.  The same
contradiction works along arbitrary sequences of maps, basepoints, and
allowed deterministic choices with their fixed common envelope.
Conditioning on independently sampled suffixes and steps, or selecting
violating deterministic choices if convergence failed, gives the stated
random-choice version.
\end{proof}
